% ==================================================
% ABSTRAK INDONESIA
% ==================================================

\clearpage

\phantomsection
\addcontentsline{toc}{chapter}{ABSTRAK}

\renewcommand{\contentsname}{ABSTRAK}

%% ===== JUDUL ABSTRAK =====
\begin{center}
    {\fontsize{14pt}{10pt}\selectfont\bfseries
    \MakeUppercase{\varJudulID}}
\end{center}

\vspace{0.2cm}

%% ===== DATA MAHASISWA =====
\hspace*{0.2cm}
\begin{tabular}{@{}p{3.9cm}@{}p{0.4cm}@{}p{10cm}@{}}
    Nama Mahasiswa & : & \varMhsNama \\
    NRP Mahasiswa & : & \varMhsNrp \\
    Pembimbing  & : & \varPromotorSatu \\
    Ko-Pembimbing & : & \varPromotorDua \\
\end{tabular}

\vspace{0.2cm}

\begin{center}
    {\fontsize{14pt}{16pt}\selectfont\bfseries
    ABSTRAK\par}
\end{center}

\vspace{0.2cm}

Isi abstrak dimulai di sini... Tuliskan abstrak proposal disertasi dalam Bahasa Indonesia pada bagian ini. Isi abstrak dapat mencakup latar belakang singkat, permasalahan penelitian, tujuan penelitian, metode yang diusulkan, serta kontribusi yang diharapkan.

Isi abstrak dimulai di sini... Tuliskan abstrak proposal disertasi dalam Bahasa Indonesia pada bagian ini. Isi abstrak dapat mencakup latar belakang singkat, permasalahan penelitian, tujuan penelitian, metode yang diusulkan, serta kontribusi yang diharapkan.

Isi abstrak dimulai di sini... Tuliskan abstrak proposal disertasi dalam Bahasa Indonesia pada bagian ini. Isi abstrak dapat mencakup latar belakang singkat, permasalahan penelitian, tujuan penelitian, metode yang diusulkan, serta kontribusi yang diharapkan.

Isi abstrak dimulai di sini... Tuliskan abstrak proposal disertasi dalam Bahasa Indonesia pada bagian ini. Isi abstrak dapat mencakup latar belakang singkat, permasalahan penelitian, tujuan penelitian, metode yang diusulkan, serta kontribusi yang diharapkan.

\vspace{0.8cm}
\noindent\textbf{Kata kunci:} semantic segmentation, deep learning, penyakit daun, lightweight model, robustness.

\clearpage

